% TEMPLATE-MILC-v3.tex - plantilla maestra UNICA de las guias MILC v3.
%
% La usan las tres familias de guias, cada una con su driver:
%   · Grados 6-11  -> scripts/build-guias-g11.py
%   · Bebras       -> scripts/build-guias-bebras.py
%   · Semillero    -> scripts/build-guias-semillero.py
%
% Antes habia tres copias casi identicas (~400 lineas iguales, 6 distintas):
% cada mejora habia que aplicarla tres veces, y en la practica no se hacia
% (el motor de recursos vivia solo en dos de ellas). Lo que varia entre
% familias esta parametrizado en 6 placeholders que cada driver llena
% (se nombran aqui SIN su sintaxis real: el builder reemplaza por texto plano,
% tambien dentro de los comentarios, y un valor multilinea rompe el preambulo):
%
%   HEADER_IZQ        encabezado izquierdo de pagina
%   HEADER_DER        encabezado derecho de pagina
%   PORTADA_ETIQUETA  rotulo sobre el numero grande (GRADO/MOMENTO/MODULO)
%   PORTADA_SERIE     linea de serie de la portada
%   PORTADA_ID        identificador de la guia en la portada
%   PORTADA_CIRCULO   el circulito de grado (°), o vacio si no aplica
%
% El resto de claves las documenta content/guias/_SCHEMA.md. El script valida
% que no queden placeholders sin reemplazar antes de escribir el archivo final.

\documentclass[11pt,letterpaper]{article}
\usepackage[margin=1.75cm]{geometry}
\usepackage[table]{xcolor}
\usepackage{fontspec}
\usepackage[spanish]{babel}
\usepackage{tikz}
% Librerías TikZ para los diagramas de las guías (bloque `recursos:`):
%  · babel  → babel en español vuelve activos caracteres como «>», y TikZ los usa
%             en sus opciones (>=stealth, ->). Sin esta librería, cualquier
%             diagrama con flechas revienta con «\language@active@arg> has an
%             extra }». Debe ir ANTES de que se dibuje nada.
%  · positioning        → sintaxis «below left=2pt and 3pt».
%  · decorations.pathreplacing → llaves ({brace}) para señalar rangos.
\usetikzlibrary{babel,positioning,decorations.pathreplacing}
\usepackage{graphicx}
% Circuitos: el trabajo de electrónica se hace hoy con micro:bit y sus sensores
% integrados (MakeCode, sin cableado), así que no se necesitan esquemáticos.
% Si algún día entran montajes con protoboard, basta descomentar esta línea:
% los diagramas .tex/.tikz declarados en `recursos:` ya se incrustan solos.
% \usepackage{circuitikz}
\usepackage[most]{tcolorbox}
\usepackage{tabularx}
\usepackage{array}
\usepackage{fancyhdr}
\usepackage{titlesec}
\usepackage{ragged2e}
\usepackage{enumitem}
\usepackage{microtype}

% Helvetica Neue no trae la flecha U+2192, asi que cada «→» escrito literalmente
% en el YAML se compilaba a NADA, en silencio (462 flechas en 61 guias: rutas del
% tipo «paleta → tipografia → lockup» salian rotas). Mapeamos el caracter a su
% version matematica, que si tiene glifo. Asi el autor puede seguir escribiendo
% «→» comodamente en el YAML.
\usepackage{newunicodechar}
\newunicodechar{→}{\ensuremath{\rightarrow}}

\setmainfont{Helvetica Neue}
\setsansfont{Helvetica Neue}
\setlength{\parindent}{0pt}
\setlength{\parskip}{6pt}
\hyphenpenalty=200
\exhyphenpenalty=200
\pretolerance=400
\tolerance=2000
\setlength{\emergencystretch}{1em}
\renewcommand{\arraystretch}{1.25}
\setlist{leftmargin=5mm,itemsep=1mm,topsep=1mm}
\newcolumntype{Y}{>{\RaggedRight\arraybackslash}X}

\definecolor{milcMagenta}{HTML}{E6007E}
\definecolor{milcVino}{HTML}{5A0038}
\definecolor{milcTurquesa}{HTML}{0093A5}
\definecolor{milcVerde}{HTML}{58A923}
\definecolor{milcMostaza}{HTML}{E5B400}
\definecolor{milcOcre}{HTML}{B86600}
\definecolor{milcNegro}{HTML}{241816}
\definecolor{milcGris}{HTML}{F7F7F4}
\definecolor{milcLinea}{HTML}{E8E2DD}
\definecolor{milcGradePrimary}{HTML}{84CC16}
\definecolor{milcGradeDark}{HTML}{4D7C0F}
\definecolor{milcGradeSoft}{HTML}{ECFCCB}
\definecolor{milcGradeAccent}{HTML}{CFE7FF}
\definecolor{milcGradeText}{HTML}{FFFFFF}
\color{milcNegro}

\pagestyle{fancy}
\fancyhf{}
\lhead{\small Tecnología e Informática · Grado 7}
\rhead{\small Guía 17 · MILC v3}
\cfoot{\small\thepage}
\renewcommand{\headrulewidth}{0pt}

\newtcolorbox{titlebox}[1]{
  enhanced,colback=#1,colframe=#1,arc=7mm,boxrule=0pt,
  left=7mm,right=7mm,top=3mm,bottom=3mm,
  before skip=8pt,after skip=8pt,
  fontupper=\sffamily\bfseries\Large\color{white}
}

\newtcolorbox{softbox}[3]{
  enhanced,colback=#2,colframe=#1,borderline west={4pt}{0pt}{#1},
  arc=2mm,boxrule=.4pt,left=4mm,right=4mm,top=3mm,bottom=3mm,
  before skip=6pt,after skip=8pt,title={#3},coltitle=white,
  colbacktitle=#1,fonttitle=\sffamily\bfseries,fontupper=\RaggedRight,
  attach boxed title to top left={xshift=3mm,yshift=-2mm},
  boxed title style={arc=2mm, boxrule=0pt}
}

\newtcolorbox{drawbox}[2]{
  enhanced,colback=white,colframe=#1,arc=4mm,boxrule=1pt,
  left=5mm,right=5mm,top=4mm,bottom=4mm,before skip=8pt,after skip=8pt,
  title={#2},coltitle=white,colbacktitle=#1,fonttitle=\sffamily\bfseries,
  fontupper=\RaggedRight,
  attach boxed title to top left={xshift=4mm,yshift=-2mm},
  boxed title style={arc=3mm, boxrule=0pt}
}

\newtcolorbox{infoband}[1]{
  enhanced,colback=#1!8,colframe=#1!50,arc=2mm,boxrule=.5pt,
  left=4mm,right=4mm,top=2mm,bottom=2mm,before skip=4pt,after skip=4pt,
  fontupper=\small\RaggedRight
}

% Puente narrativo entre secciones (contrato editorial Conéctate).
% Da continuidad: "Acabas de X. Ahora vas a Y porque Z."
% Visualmente: banda izquierda en color del grado, texto en itálica pequeña.
\newtcolorbox{puentebox}{
  enhanced,colback=milcGradeSoft,colframe=milcGradeDark,
  borderline west={3pt}{0pt}{milcGradeDark},
  arc=0mm,boxrule=0pt,left=5mm,right=4mm,top=2.5mm,bottom=2.5mm,
  before skip=4pt,after skip=8pt,
  fontupper=\itshape\small\color{milcNegro}\RaggedRight
}
% La flecha va en modo matematico a proposito: Helvetica Neue NO tiene el glifo
% U+2192, asi que \textrightarrow se compilaba a NADA («Missing character: There
% is no → in font Helvetica Neue») y el puente salia sin flecha. En modo
% matematico la flecha viene de la fuente matematica y si se dibuja.
\newcommand{\puente}[1]{%
  \begin{puentebox}%
    \textcolor{milcGradeDark}{$\rightarrow$}\hspace{2mm}#1%
  \end{puentebox}%
}

\newcommand{\linead}[1]{\noindent\color{milcLinea}\rule{#1}{0.5pt}\color{milcNegro}}
\newcommand{\respuesta}[1]{\par\vspace{2mm}\foreach \n in {1,...,#1}{\noindent\color{milcLinea}\rule{\linewidth}{0.5pt}\par\vspace{5mm}}\color{milcNegro}}
\newcommand{\checkbox}{\textcolor{milcGradeDark}{\fbox{\phantom{x}}}\hspace{5pt}}

\newcommand{\phasecard}[4]{
\begin{tcolorbox}[enhanced,colback=#3,colframe=#2,arc=3mm,boxrule=.5pt,height=3.05cm,valign=center,halign=center,left=1.5mm,right=1.5mm,top=2mm,bottom=2mm]
{\sffamily\bfseries\fontsize{22}{24}\selectfont\textcolor{#2}{#1}}\par
{\sffamily\bfseries\small #4}
\end{tcolorbox}
}

% ── Recursos visuales de la guía ──────────────────────────────────────
% Declarados en el bloque `recursos:` del YAML y emitidos por el builder.
% \guiaFigura   → imagen rasterizada o PDF (foto, carta celeste, montaje).
% \guiaDiagrama → fuente TikZ/circuitikz (.tex) incrustada como vector:
%                 es la vía para circuitos y esquemáticos editables.
% #1 = ruta relativa al .tex · #2 = pie (puede ir vacío)
\newcommand{\guiaFigura}[2]{%
\begin{center}
\includegraphics[width=0.88\linewidth,height=0.38\textheight,keepaspectratio]{#1}\\[1.5mm]
{\footnotesize\itshape\textcolor{milcNegro!75}{#2}}
\end{center}
}

\newcommand{\guiaDiagrama}[2]{%
\begin{center}
\input{#1}\\[1.5mm]
{\footnotesize\itshape\textcolor{milcNegro!75}{#2}}
\end{center}
}

\begin{document}
\thispagestyle{empty}
\begin{tikzpicture}[remember picture,overlay]
  \fill[white] (current page.south west) rectangle (current page.north east);
  \path[fill=milcGradePrimary,rounded corners=16mm]
    ([xshift=.55cm,yshift=-.55cm]current page.north west) rectangle
    ([xshift=-.55cm,yshift=3.65cm]current page.south east);

  \node[anchor=north west,text=milcGradeText] at ([xshift=2.0cm,yshift=-1.85cm]current page.north west)
    {\sffamily\bfseries\fontsize{14}{16}\selectfont GRADO};
  \node[anchor=north west,text=milcGradeText,opacity=.82] at ([xshift=2.0cm,yshift=-2.75cm]current page.north west)
    {\sffamily\bfseries\fontsize{9}{11}\selectfont SERIE GUÍAS · TECNOLOGÍA E INFORMÁTICA};

  \begin{scope}[shift={([xshift=-3.05cm,yshift=-2.55cm]current page.north east)}]
    \fill[white,opacity=.80] (-.62,-.52) rectangle (.62,.52);
    \draw[milcGradeText,line width=1pt] (-.62,-.52) rectangle (.62,.52);
    \fill[milcGradeDark] (-.42,-.38) rectangle (-.22,.06);
    \fill[milcGradeAccent] (-.10,-.38) rectangle (.10,.24);
    \fill[milcGradePrimary!60!black] (.22,-.38) rectangle (.42,.38);
  \end{scope}

  \node[anchor=west,text=milcGradeText] at ([xshift=1.9cm,yshift=-12.85cm]current page.north west)
    {\sffamily\bfseries\fontsize{124}{124}\selectfont 7};
  % El circulito de grado (°) solo aplica a las guias de grado («11°»); en
  % Bebras y Semillero el numero grande es un momento/modulo, no un grado.
  \draw[milcGradeText,line width=1.7pt]
    ([xshift=4.52cm,yshift=-11.68cm]current page.north west) circle (.13cm);

  \node[anchor=west,text=milcGradeText,text width=8.7cm,align=left] at ([xshift=10.35cm,yshift=-11.6cm]current page.north west)
    {
     {\sffamily\bfseries\fontsize{19}{22}\selectfont Bucles ---\\repeticiones\\inteligentes\\en los\\algoritmos}\\[4mm]
     {\sffamily\bfseries\fontsize{23}{26}\selectfont Guía 17-7-TIC}\\[2mm]
     {\sffamily\fontsize{13}{16}\selectfont Período 2 · Algoritmia y pensamiento computacional}\\[5mm]
     {\sffamily\bfseries\fontsize{11}{13}\selectfont Institución Educativa Sor María Juliana}\\[1mm]
     {\sffamily\fontsize{10}{12}\selectfont Autor: Dr. Álvaro Cárdenas Orozco}
    };

  \path[fill=milcGradeSoft,rounded corners=7mm]
    ([xshift=1.35cm,yshift=.85cm]current page.south west) rectangle
    ([xshift=-1.35cm,yshift=4.25cm]current page.south east);
  \draw[milcGradeDark,line width=.65pt,rounded corners=7mm]
    ([xshift=1.35cm,yshift=.85cm]current page.south west) rectangle
    ([xshift=-1.35cm,yshift=4.25cm]current page.south east);
  \node[anchor=center,text=milcNegro,text width=17.0cm,align=center] at ([yshift=2.55cm]current page.south)
    {\sffamily\fontsize{10}{12}\selectfont Metodología MILC · Apertura ancestral · Pensamiento computacional · Triángulo Dussel-Estoicismo-Floridi\\
    \textcolor{milcGradeDark}{\bfseries Producto:} Un \textbf{algoritmo que use al menos 2 bucles} (uno de tipo \emph{repetir N veces}, otro de tipo \emph{repetir hasta que}). Tema sugerido: \emph{``Tejer una pulsera de 30 vueltas''}, \emph{``Saludar a los 20 invitados a una fiesta''}, \emph{``Lavarse los dientes (cepillar arriba, abajo, izquierda, derecha)''}. Más \textbf{su diagrama de flujo} (con el bucle dibujado correctamente). Más \textbf{tabla comparativa} de los 2 tipos de bucles en el cuaderno.};
\end{tikzpicture}
\mbox{}
\newpage

\section*{Apertura: Bucles --- repeticiones inteligentes en los algoritmos}

\begin{softbox}{milcGradeDark}{milcGradeSoft}{Saber ancestral}
\textbf{Doña Sofía la tejedora del Valle del Cauca repetía el mismo patrón 200 veces para tejer una mochila.} En la vereda La Plata, las tejedoras de canastos del Cauca, las hamaqueras del Pacífico, las costureras de Buga tenían un saber compartido: \emph{el patrón se repite}. Cuando doña Sofía empezaba una mochila, sabía que tenía que ejecutar la secuencia \emph{``levantar hebra, pasarla sobre 2, pasarla por debajo de 1''} \textbf{exactamente 200 veces}. \emph{¿Contaba mentalmente?} Sí. \emph{¿Se aburría?} No --- la repetición rítmica era \emph{meditativa}. \emph{¿Se equivocaba?} A veces, y entonces destejía hasta el último punto correcto y volvía a empezar. Esa estructura --- \emph{``repetir N veces''} o \emph{``repetir hasta lograr X''} --- es exactamente lo que en programación llamamos \textbf{bucle} o \textbf{ciclo}. Sin bucles, una mochila Wayuu tomaría escribir 200 instrucciones separadas. \textbf{Con bucles, una instrucción se ejecuta automáticamente N veces}. Es la elegancia algorítmica que las tejedoras ya conocían.
\end{softbox}

\begin{softbox}{milcTurquesa}{milcGris}{Saber contemporáneo}
Un \textbf{bucle} (o \emph{ciclo}, o \emph{loop} en inglés) es una estructura que \textbf{repite una o varias acciones hasta que se cumple una condición}. Es el tercer componente fundamental de la programación (después de variables y condicionales). \textbf{Hay 2 tipos básicos.} \textbf{(1) Bucle ``repetir N veces''} (\emph{for}): se repite un número fijo de veces. Ejemplo: \emph{REPETIR 10 VECES: dar pasos}. Útil cuando sabes exactamente cuántas repeticiones quieres. \textbf{(2) Bucle ``repetir hasta que''} (\emph{while}): se repite mientras una condición es verdadera. Ejemplo: \emph{REPETIR HASTA QUE peso = 100kg: comer}. Útil cuando no sabes cuántas repeticiones, pero sabes cuándo parar. \textbf{Los bucles son lo que hace que un programa pueda procesar muchos datos sin escribir mil veces lo mismo}: leer todos los correos, mover todos los archivos, sumar todas las notas. Sin bucles, los programas serían imposibles de escribir.
\end{softbox}

\begin{softbox}{milcMostaza}{milcGris}{Pregunta-puente}
\textit{Si tuvieras que escribir un programa para \emph{``saludar a los 30 estudiantes de la clase, uno por uno, llamándolos por su nombre''}, ¿\textbf{escribirías 30 instrucciones separadas} o usarías un bucle?}
\end{softbox}

\begin{softbox}{milcVerde}{milcGris}{Saber y saber hacer}
Al terminar la clase: (1) podrás \textbf{identificar} los 2 tipos de bucles; (2) sabrás \textbf{aplicar} cada tipo a su situación; (3) podrás \textbf{evaluar} si un algoritmo necesita bucle o no; (4) habrás \textbf{creado} un algoritmo con 2 bucles + diagrama.
\end{softbox}



\small
\begin{tabularx}{\textwidth}{>{\bfseries}p{3.0cm}Y}
\rowcolor{milcGradePrimary!12}Autor & Dr. Álvaro Cárdenas Orozco\\
DBA & Construye algoritmos y diagramas de flujo para resolver problemas paso a paso (MEN, T\&I 7°)\\
Referentes & Saber ancestral del tejedor · Pensamiento computacional · Scratch\\
Producto final & Un \textbf{algoritmo que use al menos 2 bucles} (uno de tipo \emph{repetir N veces}, otro de tipo \emph{repetir hasta que}). Tema sugerido: \emph{``Tejer una pulsera de 30 vueltas''}, \emph{``Saludar a los 20 invitados a una fiesta''}, \emph{``Lavarse los dientes (cepillar arriba, abajo, izquierda, derecha)''}. Más \textbf{su diagrama de flujo} (con el bucle dibujado correctamente). Más \textbf{tabla comparativa} de los 2 tipos de bucles en el cuaderno.\\
\end{tabularx}
\normalsize

\puente{Hoy haces 4 cosas: identificas repeticiones en situaciones reales, aprendes los 2 tipos de bucles, escribes algoritmo con 2 bucles, dibujas el diagrama.}

\begin{titlebox}{milcGradeDark}
Ruta MILC: así vamos a aprender
\end{titlebox}
\begin{tabularx}{\textwidth}{>{\centering\arraybackslash}X>{\centering\arraybackslash}X>{\centering\arraybackslash}X>{\centering\arraybackslash}X}
\phasecard{1}{milcVerde}{milcGris}{Escucha\\[-1mm]\small Observo, escucho y pregunto.} &
\phasecard{2}{milcTurquesa}{milcGris}{Sistematización\\[-1mm]\small Pensamiento computacional.} &
\phasecard{3}{milcMagenta}{milcGris}{Praxis\\[-1mm]\small Construyo y aplico.} &
\phasecard{4}{milcVino}{milcGris}{Evaluación\\[-1mm]\small Reflexión liberadora.}
\end{tabularx}


\puente{Empezamos viendo qué pasaría sin bucles (caos).}

\begin{titlebox}{milcVerde}
Fase 1 · Escucha
\end{titlebox}

\begin{softbox}{milcVerde}{milcGris}{Escena para escuchar}
\textbf{Actividad 1 · ANALIZA --- ¿Cuántas instrucciones sin bucle?} (10 min · individual). Imagina este problema: \emph{tienes que dar la bienvenida a los 30 estudiantes nuevos de 7°. A cada uno: ``Hola, bienvenido a la clase''.} \textbf{En el cuaderno responde:} \emph{(1) Sin usar bucles, ¿cuántas instrucciones escribirías?}. \emph{(2) ¿Cuánto te tomaría escribirlas?}. \emph{(3) ¿Qué pasa si en lugar de 30 son 300 estudiantes?}. \emph{(4) ¿Crees que hay una manera más inteligente?}.
\end{softbox}

\checkbox Pensé en el problema sin bucle.\\
\checkbox Calculé cuántas instrucciones serían.\\
\checkbox Reconocí que escalar es problema sin bucles.

\begin{infoband}{milcVerde}


\textbf{Lo que va al cuaderno (Actividad 1 · ANALIZA).} Título: «Actividad 1 --- 30 saludos sin bucle». \textbf{Formato:} 4 respuestas cortas. \textbf{Extensión:} 5-6 líneas. \textbf{Sabes que terminaste cuando} ves que sin bucles los programas serían inmanejables.
\end{infoband}

\puente{Después aprendes los 2 tipos + el contador como variable amiga.}

\begin{titlebox}{milcTurquesa}
Fase 2 · Sistematización con pensamiento computacional
\end{titlebox}

\begin{softbox}{milcTurquesa}{milcGris}{Cuatro pilares aplicados al tema}
La intuición es correcta: \textbf{sin bucles, los programas se vuelven imposibles}. Ahora aprende los 2 tipos de bucles y cómo se escriben.
\end{softbox}

\small
\begin{tabularx}{\textwidth}{>{\bfseries\RaggedRight\arraybackslash}p{3.6cm}Y}
\rowcolor{milcTurquesa!90}\color{white}Pilar & \color{white}\bfseries Aplicación al tema\\
1. Descomposición & \textbf{Bucle tipo 1 --- ``Repetir N veces'' (for).} Cuando \emph{sabes exactamente cuántas veces} quieres repetir. Estructura: \emph{REPETIR [N] VECES: [acción]}. Ejemplo del saludo: \emph{REPETIR 30 VECES: decir ``Hola, bienvenido''}. Eso es 1 instrucción que ejecuta 30 veces. En programación a veces se usa una \textbf{variable contador} que va de 1 a N: \emph{PARA cada estudiante DESDE 1 HASTA 30: decir ``Hola estudiante [contador]''}. La variable contador toma valores 1, 2, 3, ..., 30 sucesivamente.\\
2. Reconocimiento de patrones & \textbf{Bucle tipo 2 --- ``Repetir hasta que'' (while).} Cuando \emph{no sabes cuántas veces} pero sabes la condición de parada. Estructura: \emph{REPETIR HASTA QUE [condición]: [acción]}. Ejemplo: \emph{REPETIR HASTA QUE puntaje $\geq$ 100: jugar un nivel}. Se repite mientras puntaje sea menor que 100. Cuando puntaje llega a 100 (o más), el bucle termina. Otra variante: \emph{MIENTRAS [condición es verdadera]: [acción]} --- igual de común, lógica equivalente.\\
3. Abstracción & \textbf{El contador como variable amiga del bucle.} Casi todo bucle usa una \textbf{variable contador}: número que se incrementa cada vuelta del bucle. Ejemplo de tejido: \texttt{vueltaActual = 0}; \emph{REPETIR HASTA QUE vueltaActual = 200: tejer una vuelta; vueltaActual = vueltaActual + 1}. Cada iteración: teje 1 vuelta y suma 1 al contador. Cuando llega a 200, termina. El contador es lo que conecta variables (S5) con bucles (S7).\\
4. Algoritmo conceptual & \textbf{El bucle infinito (cuidado).} Un bucle que NUNCA termina porque la condición nunca se vuelve falsa. Ejemplo MAL: \emph{REPETIR HASTA QUE 5 > 10: jugar} --- esto es siempre falso, el bucle no para nunca, el programa se congela. \textbf{Cómo evitarlo:} (a) verificar que la condición se pueda volver verdadera en algún momento; (b) verificar que algo cambia adentro del bucle (incrementar contador, modificar variable). En Scratch y otros lenguajes, los bucles infinitos hacen que el programa se trabe.\\
\end{tabularx}
\normalsize

\begin{drawbox}{milcTurquesa}{2 tipos de bucles + contador + cuidado infinito}
\textbf{Tipo 1 (for):} REPETIR [N] VECES: [acción]. \\[1mm]
\textbf{Tipo 2 (while):} REPETIR HASTA QUE [condición]: [acción]. \\[1mm]
\textbf{Contador:} variable que cuenta las iteraciones (vueltaActual = 0; vueltaActual = vueltaActual + 1). \\[1mm]
\textbf{Cuidado bucle infinito:} verificar que la condición se pueda cumplir.
\end{drawbox}

\begin{softbox}{milcMostaza}{milcGris}{Errores comunes y su corrección}
\textbf{Errores frecuentes con bucles.} (1) \emph{Olvidar incrementar el contador} --- bucle infinito porque nunca cambia la condición. (2) \emph{Condición que nunca se vuelve verdadera} --- bucle infinito. \emph{``REPETIR HASTA QUE 0 = 1''} no termina nunca. (3) \emph{Usar bucle cuando no se necesita} --- si solo es 1 vez, no necesitas bucle. (4) \emph{Usar tipo for cuando debería ser while} (o viceversa) --- depende de si conoces el número exacto o solo la condición. (5) \emph{Anidar bucles sin necesidad} --- 5 bucles dentro de bucles se vuelven inentendibles. Mejor descomponer.
\end{softbox}

\begin{infoband}{milcTurquesa}


\textbf{Lo que va al cuaderno (Actividad 2 · EXPLICA).} Título: «Actividad 2 --- 2 tipos de bucles». \textbf{Formato:} tabla comparativa de 4 filas, 3 columnas (criterio, FOR, WHILE). Filas: \emph{cuándo se usa}, \emph{ejemplo de la vida real}, \emph{estructura escrita}, \emph{cómo se ve en diagrama}. \textbf{Extensión:} 4 filas. \textbf{Sabes que terminaste cuando} podrías escoger el tipo correcto de bucle para cualquier situación.
\end{infoband}

\puente{Con todo claro, escribes tu algoritmo con bucles + diagrama.}

\begin{titlebox}{milcMagenta}
Fase 3 · Praxis con pensamiento computacional
\end{titlebox}

\begin{softbox}{milcMagenta}{milcGris}{Construyendo el producto con los cuatro pilares}
\textbf{Actividad 3 · APLICA --- Algoritmo con 2 bucles + diagrama} (30 min · individual). Vas a escribir un algoritmo que use \textbf{2 bucles distintos} (uno de cada tipo) y dibujarlo en diagrama de flujo.
\end{softbox}

\small
\begin{tabularx}{\textwidth}{>{\bfseries\RaggedRight\arraybackslash}p{3.6cm}Y}
\rowcolor{milcMagenta!90}\color{white}Pilar & \color{white}\bfseries Aplicación al producto\\
Descomposición & \textbf{Paso 1 --- Escoge el tema del algoritmo.} Tiene que ser algo que requiera \emph{2 tipos de repetición}: una \emph{N veces} (sabes cuántas) y una \emph{hasta que} (no sabes cuántas, sabes cuándo parar). Sugerencias: \emph{``Tejer una pulsera de 30 vueltas pero hasta que se acabe el hilo''} (for por las 30 vueltas, while por el hilo). \emph{``Saludar a los 20 invitados de una fiesta y hasta que se acaben los bocaditos''} (for por los 20, while por los bocaditos). \emph{``Estudiar 5 ejercicios pero hasta que entienda el tema''}.\\
Patrones & \textbf{Paso 2 --- Identifica las variables y el contador.} Lista las variables: \texttt{cantidadVueltas} (número), \texttt{hiloRestante} (número), \texttt{contadorVueltas} (número, empieza en 0). Identifica cuál es el contador para cada bucle.\\
Abstracción & \textbf{Paso 3 --- Escribe el algoritmo.} Estructura sugerida: \emph{INICIO. Variables: cantidadVueltas = 30, hiloRestante = 100, contadorVueltas = 0. BUCLE 1 (for): PARA contadorVueltas DESDE 1 HASTA 30: tejer 1 vuelta; hiloRestante = hiloRestante - 1. BUCLE 2 (while): MIENTRAS hiloRestante > 0: tejer 1 vuelta más; hiloRestante = hiloRestante - 1. FIN.}. Adapta a tu tema.\\
Algoritmo de armado & \textbf{Paso 4 --- Dibuja el diagrama de flujo.} Los bucles se representan así en diagrama: \textbf{(a)} óvalo INICIO. \textbf{(b)} Para el bucle ``repetir N veces'': rombo \emph{``¿contador < N?''}, flecha SÍ va al rectángulo de la acción + incremento del contador, flecha NO va al siguiente paso. La acción y el incremento vuelven al rombo (creando el ciclo visual). \textbf{(c)} Para el bucle ``repetir hasta que'': rombo \emph{``¿condición?''}, flecha SÍ va al rectángulo de la acción y vuelve al rombo, flecha NO sale del bucle. \textbf{(d)} óvalo FIN.\\
\end{tabularx}
\normalsize

\begin{drawbox}{milcMagenta}{Antes de cerrar la S7}
\checkbox Mi algoritmo tiene 2 bucles, uno de cada tipo. \\[1mm]
\checkbox El contador del bucle for está identificado. \\[1mm]
\checkbox El bucle while tiene condición que se vuelve falsa eventualmente. \\[1mm]
\checkbox El diagrama de flujo muestra los 2 bucles visualmente. \\[1mm]
\checkbox La tabla comparativa for vs while está completa. \\[1mm]
\checkbox Verifiqué que NO hay bucle infinito en mi algoritmo.
\end{drawbox}

\begin{softbox}{milcMagenta}{milcGris}{Plantilla de trabajo}
\textbf{Estructura.} \textit{Paso 1:} tema del algoritmo (que requiera 2 tipos de repetición). \textit{Paso 2:} variables + contador. \textit{Paso 3:} algoritmo con bucle for + bucle while. \textit{Paso 4:} diagrama de flujo. \textit{Paso 5:} tabla comparativa. \textit{Paso 6:} verificación + reflexión + firma.
\end{softbox}

\begin{infoband}{milcMagenta}


\textbf{Lo que va al cuaderno (Actividad 3 · APLICA).} Título: «Actividad 3 --- Algoritmo con 2 bucles». \textbf{Formato:} variables + algoritmo + diagrama + tabla + verificación. \textbf{Extensión:} 1 página y media. \textbf{Sabes que terminaste cuando} podrías ejecutar mentalmente tu algoritmo y verificar que termina (no es infinito).
\end{infoband}

\puente{El producto es algoritmo + diagrama + tabla comparativa.}

\begin{titlebox}{milcMostaza}
Producto de la sesión
\end{titlebox}
\textbf{Algoritmo con 2 bucles + diagrama · for y while aplicados · sin bucle infinito}.
\begin{softbox}{milcMostaza}{milcGris}{Criterios de calidad}
(1) El algoritmo tiene 1 bucle for (repetir N veces) bien aplicado. (2) El algoritmo tiene 1 bucle while (repetir hasta que) bien aplicado. (3) El diagrama de flujo muestra los 2 bucles visualmente. (4) El contador del bucle for está identificado. (5) Se verificó que ningún bucle es infinito.
\end{softbox}

\puente{Antes de salir, verifica que tu algoritmo NO entre en bucle infinito (que sí termine).}

\begin{titlebox}{milcVino}
Fase 4 · Evaluación liberadora — cinco dimensiones
\end{titlebox}

\begin{softbox}{milcVino}{milcGris}{Sentido de la evaluación}
Evaluar no es solo revisar una nota. Es reconocer qué aprendí, qué cambió en mí, qué emociones manejé, cómo me relaciono con la ciudad y la comunidad, y qué le dejaré a quien viene detrás.
\end{softbox}

\textbf{Mi reflexión liberadora en cinco dimensiones.} Completa cada frase iniciada:

\small
\begin{tabularx}{\textwidth}{>{\bfseries\RaggedRight\arraybackslash}p{3.4cm}Y>{\RaggedRight\arraybackslash}p{4.6cm}}
\rowcolor{milcVino}\color{white}Dimensión & \color{white}\bfseries Mi respuesta (frame starter) & \color{white}\bfseries Algo que descubrí\\
1. Personal & Lo que más cambió en mí fue \linead{6.5cm} & \linead{4.2cm}\\
2. Emocional & La emoción más fuerte fue \linead{2.5cm} y la regulé \linead{3cm} & \linead{4.2cm}\\
3. Ciudadana & En democracia esto importa porque \linead{6cm} & \linead{4.2cm}\\
4. Local (Cartago) & En mi barrio sería útil para \linead{6.5cm} & \linead{4.2cm}\\
5. Intergeneracional & A mi abuela/abuelo le diría \linead{6.5cm} & \linead{4.2cm}\\
\end{tabularx}
\normalsize

\begin{infoband}{milcVino}
\textbf{Para la próxima sustentación.} Vuelve a esta tabla en seis meses. Verás que las respuestas cambian — no porque lo aprendido cambie, sino porque tú cambias. La evaluación liberadora es una conversación contigo a través del tiempo.
\end{infoband}


\puente{Cerramos con 3 voces sobre por qué saber repetir bien es virtud antigua.}

\begin{titlebox}{milcGradeDark}
Triángulo de pensamiento · cierre filosófico
\end{titlebox}

\begin{softbox}{milcVino}{milcGris}{Dussel · lente del nosotros}
\textit{````La repetición no es aburrida cuando entiendes su sentido. La tejedora tejió 200 vueltas porque cada vuelta importaba.''''}

Hay una idea moderna de que la \emph{repetición es aburrida} y debe automatizarse. Es cierto a medias: \textbf{lo que el computador puede repetir automáticamente, no debes repetirlo a mano}. Pero \emph{algunas repeticiones son meditativas y formativas}: doña Sofía tejiendo, el atleta practicando, el músico ensayando. \emph{Los bucles automatizan las repeticiones mecánicas para que tú te concentres en las repeticiones formativas}.

\textbf{¿Qué repeticiones de mi vida son mecánicas (deberían automatizarse) y cuáles son formativas (vale la pena repetir conscientemente)?}
\end{softbox}
\linead{\linewidth}\\[1mm]
\linead{\linewidth}

\begin{softbox}{milcMostaza}{milcGris}{Estoicismo · Marco Aurelio · cuidado interior}
\textit{````Lo que se repite con criterio, perfecciona. Lo que se repite por inercia, embrutece. La diferencia está en la atención.''''}

\textbf{Los bucles bien diseñados son como la disciplina diaria}: hacen lo mismo cada vez, pero con propósito claro y condición de parada. Sin propósito (bucle infinito), repetir embrutece. Con propósito (condición clara), repetir perfecciona. Marco Aurelio escribía sus reflexiones cada noche --- bucle con propósito, no inercia.

\textbf{Las repeticiones de mi vida diaria, ¿tienen propósito claro o las hago por inercia?}
\end{softbox}
\linead{\linewidth}\\[1mm]
\linead{\linewidth}

\begin{softbox}{milcTurquesa}{milcGris}{Floridi · lente de la infoesfera}
\textit{````Los bucles son lo que hizo posible la computación masiva. Sin bucles, no habría Google ni Wikipedia.''''}

Google indexa millones de páginas \textbf{usando bucles}: \emph{REPETIR por cada página: leer, analizar, indexar, pasar a la siguiente}. Wikipedia tiene millones de artículos. Spotify procesa millones de canciones. \emph{Los bucles son lo que escaló la informática desde calculadoras simples hasta la infraestructura del mundo digital}. Entenderlos en 7° es entender una pieza fundamental del mundo moderno.

\textbf{¿Qué cosas del mundo digital que uso a diario serían imposibles sin bucles?}
\end{softbox}
\linead{\linewidth}\\[1mm]
\linead{\linewidth}

\begin{titlebox}{milcVino}
Compromiso final
\end{titlebox}

\small
\begin{tabularx}{\textwidth}{>{\RaggedRight\arraybackslash}p{3.0cm}YYY}
\rowcolor{milcVino}\color{white}\bfseries Criterio & \color{white}\bfseries Lo logré & \color{white}\bfseries En proceso & \color{white}\bfseries Necesito apoyo\\
Comprensión del tema & Explico ideas centrales con mis palabras. & Reconozco ideas pero falta precisión. & Necesito repasar conceptos base.\\
Pensamiento computacional & Apliqué los 4 pilares al diseñar mi producto. & Apliqué algunos pilares. & Necesito releer la fase 2.\\
Aplicación y producto & Mi producto cumple los criterios y se puede revisar. & Producto funcional con ajustes pendientes. & Aún en construcción.\\
Reflexión 5 dimensiones & Respondí cada dimensión con honestidad. & Respondí algunas con detalle. & Mis respuestas son muy generales.\\
Triángulo de pensamiento & Conecté las 3 voces con mi trabajo real. & Conecté al menos una. & No relaciono las voces con mi proceso.\\
\end{tabularx}
\normalsize

\textbf{Hoy entendí que:}\\[1mm]
\linead{\linewidth}\\[1mm]
\linead{\linewidth}

\textbf{Mi compromiso para la próxima sesión es:}\\[1mm]
\linead{\linewidth}\\[1mm]
\linead{\linewidth}

\begin{softbox}{milcMostaza}{milcGris}{Cita de despedida · Marco Aurelio}
\textit{``El obstáculo en el camino se vuelve el camino. Lo que estorba al camino se vuelve camino.''}\\[1mm]
Cada error de hoy, bien observado, se convierte en pieza del camino que pisarás mañana.
\end{softbox}

\small
\begin{tabularx}{\textwidth}{>{\bfseries}p{3.6cm}Y}
\rowcolor{milcGradeDark!12}Nombre completo: & \linead{10cm}\\
Fecha: & \linead{4cm} \quad Curso/grupo: \linead{4cm}\\
Mi firma: & \linead{9cm}\\
\end{tabularx}
\normalsize

\begin{infoband}{milcGradeDark}
\textbf{Para la próxima sesión.} Esta semana, identifica 3 situaciones reales donde te tocaría usar bucles si fueras a automatizarlas. La próxima clase entramos a \textbf{Scratch}: abres tu primer programa con bloques. Vas a programar de verdad.
\end{infoband}

\end{document}
